\section{Introduction}

Robots are increasingly capable of replacing routine manual and cognitive tasks, there remains a persistent need of human cognitive skills that machines cannot yet replicate. This technological shift is occurring against a backdrop of severe labor shortages, largely driven by the fact that many humans are unwilling to perform roles perceived as physically demanding  that will be replacible by the robots in the future.  

Industries such as agriculture, forestry, and earthmoving (Heavy-Duty Mobile Machinery) are currently facing severe skilled labor shortages and high turnover rates. This is largely driven by the harsh work environments and poor safety records associated with these roles, making them unattractive to the modern workforce

Research on adolescents confirms that human expectations are directly influenced by the prospect of automation. When students gain immediate knowledge that a specific job has a high likelihood of being replaced by AI, their willingness to choose that career significantly decreases

Automation is progressively reducing reliance on human labor in technical domains. However, current robotic systems remain limited in handling complex, dynamic tasks, resulting in a continued demand for skilled human operators. At the same time, the acquisition of such skills is difficult and time-intensive, particularly due to the need for developing muscle memory through repeated practice.

In the study of learning and knowledge, near transfer refers to the application of skills or knowledge to a similar context, while far transfer refers to its application in a dissimilar context. Although the near transfer is only repeating similar jobs, it still needs a certain amount of repeating for humans to form a muscle memory.  Prior research has explored Extended Reality (XR) for training in domains such as medical procedures and Aviation and Aerospace. 

Current drone simulators are largely entertainment-oriented, featuring steep learning curves and limited support for guided skill acquisition. Unsuited controller compatibility also brings high preconfiguration frustration.

Low expectations, high pre-learning preparation, and high frustration levels during the learning process cause the beginners to lose perseverance and fail to develop muscle memory. 

The goal of this study is to design an XR-based training system that lowers the learning barrier for FPV quadcopter piloting by supporting muscle memory acquisition through AI assisted co-training and interaction desgin.

There are three research questions will be discussed in this paper:

\begin{itemize}
    \item (1) How can FPV quadcopter flight physics be simulated ?
    \item (2) How can a human-AI co-training framework be designed to assist users?
    \item (3) What interaction instruments enable high-fidelity XR prototyping?
\end{itemize}

This study adopts a design-based research approach, using agile spikes to implementing iterative prototyping of an XR FPV quadcopter simulator with physics simulation, AI co-training (PID control and Reinforcement Learning), and interaction design.

% --- Deprecated
\subsection{Background (controversial or important information, framework)}
Automation is progressively reducing reliance on human labor in technical domains. However, current robotic systems remain limited in handling complex, dynamic tasks, resulting in a continued demand for skilled human operators. At the same time, the acquisition of such skills is difficult and time-intensive, particularly due to the need for developing muscle memory through repeated practice.

Skill learning consists of two components: conceptual understanding and physical training. While conceptual knowledge can be transferred efficiently, physical skill acquisition remains a bottleneck. Extended Reality (XR) has been proposed as a medium for immersive interaction and simulation, offering potential for embodied learning.

% --- Deprecated
\subsection{Previous Studies (very short)}
Prior research has explored XR for training in domains such as medical procedures and flight simulation. AI methods, including control systems and learning-based approaches, have also been applied to assist training and automation.


\subsection{Need for Additional Research}
Existing work primarily focuses on concept transfer and high-level training outcomes, with limited attention to early-stage muscle memory development. Additionally, current drone simulators are largely entertainment-oriented, featuring steep learning curves and limited support for guided skill acquisition.

\subsection{Idea (problem)}
There is a gap between the demand for skilled operators and the accessibility of training systems. In particular, existing approaches do not effectively support the formation of muscle memory in complex physical tasks, nor do they integrate AI assistance and XR interaction in a unified framework.

\subsection{Goal}
The goal of this study is to design an XR-based training system that lowers the learning barrier for FPV quadcopter piloting by supporting muscle memory acquisition through AI-assisted co-training and interaction design.

\subsection{Research Questions}
\begin{itemize}
    \item (1) How can FPV quadcopter flight physics be simulated to support realistic and effective training?
    \item (2) How can a human-AI co-training framework be designed to assist users during skill acquisition?
    \item (3) What interaction instruments enable high-fidelity XR prototyping for physical skill training?
\end{itemize}

\subsection{Methods used (very short)}
This study adopts a design-based research approach, implementing iterative prototyping of an XR FPV drone simulator integrating physics simulation, AI co-training (PID and Reinforcement Learning), and interaction design.